\begin{answers}{ch:embedding}{埋め込みとニューラル言語モデル}
\item 番号の列 $[3,1]$ の埋め込みは、$\mE$ の 3 行目と 1 行目を並べた $(2,-1)$、$(0,1)$ です。
番号 3 の one-hot ベクトルは $\vx=(0,0,0,1)^\T$ です。
式\eqref{eq:emb-onehot}と同じく横に寝かせて掛けると、
$\vx^\T\mE=0\cdot(1,0)+0\cdot(0,1)+0\cdot(1,1)+1\cdot(2,-1)=(2,-1)$ となります。
表の 3 行目がそのまま出てきました。
表を引くことと、one-hot に行列を掛けることは、同じ計算です。

\item 合計は 1,224,960 個で、いちばん多いのは出力層（$\mW_2$、$\vb_2$）です。
式\eqref{eq:emb-concat}〜\eqref{eq:emb-logits}の部品ごとに数えます。
埋め込み表 $\mE$ は $V\times d=4096\times32=131{,}072$ 個です。
並べたベクトルの長さは $nd=160$ なので、$\mW_1$ は $256\times160=40{,}960$ 個、$\vb_1$ は 256 個です。
$\mW_2$ は $V\times d_h=4096\times256=1{,}048{,}576$ 個、$\vb_2$ は 4096 個です。
出力層だけで全体の約 86\% を占めます。
$V$ が大きいと、$V$ に比例する部品（埋め込み表と出力層）が大半になり、窓の幅 $n$ が効くのは $\mW_1$ だけです。

\item 小説では本文の $d=64$ がいちばん良く、ドリルは $d$ によらずほぼ同じでした。傾向は同じではありません。
$n=3$、$d_h=128$、更新 3000 回のまま $d$ だけを変えて実測した、検証損失は次のとおりです。
$d=16$ は小説 5.658、ドリル 2.059（パラメータ約 60 万）。
$d=64$ は小説 5.512、ドリル 2.059（本文の値）。
$d=256$ は小説 5.725、ドリル 2.055（約 168 万）でした。
小説では、$d=16$ は表現の幅が足りません。$d=256$ は訓練損失も 3.857 と $d=64$（3.787）より高く、3000 回の更新では学習しきれていないようです。
ドリルは窓の幅 3 では問題文が見えないので（\tabref{tab:emb-window}）、$d$ を変えても損失は 2.06 前後で止まります。正答率も 2\%（$d=64$ は 2.5\%）で、ほとんど変わりません。
足りないのは、ベクトルの幅ではなく、見える範囲です。

\item 1 文字目には幅 8、3 文字目には幅 10 が必要です。
「\code{123+456=}」は、数字 6 個と「\code{+}」「\code{=}」で 8 トークンです。
答えの 1 文字目を当てる時点では、窓の末尾が「\code{=}」なので、幅 8 で問題文がすべて入ります。
\tabref{tab:emb-window}で、$n=8$ のとき「窓に入る」が 100\% になったのと合っています。
3 文字目を当てる時点では、すでに書いた答えの 2 文字「\code{5}」「\code{7}」も窓に入るので、$8+2=10$ が必要です。
幅 8 のままだと、先頭の「\code{1}」「\code{2}」が窓から押し出されます。
同じ表で、$n=8$ の「最後まで」が 75\% にとどまり、$n=12$ で 100\% になるのは、このためです。
答えを書き進めるほど、問題文は窓の外へ追い出されていきます。

\item 上限は $0.5^{20}=1/1{,}048{,}576\approx9.5\times10^{-7}$、約 100 万分の 1 です。
$u_t=w\,h_{t-1}+x_t$ とおくと、連鎖律と $\tanh$ の微分 $1-\tanh^2(u)$ より
$\partial h_t/\partial h_{t-1}=(1-\tanh^2(u_t))\,w$ です。
$\tanh$ の値は $-1$ から 1 なので、$0\le1-\tanh^2(u_t)\le1$ です。
よって $\left|\partial h_t/\partial h_{t-1}\right|\le|w|$ となります。
$k$ 個前までは、連鎖律で
$\partial h_t/\partial h_{t-k}=\prod_{j=0}^{k-1}\partial h_{t-j}/\partial h_{t-j-1}$ と 1 段ずつの積になります。
各段の大きさが $|w|$ 以下なので、積の大きさは $|w|^k$ 以下です。
$|w|<1$ なら、さかのぼるほど指数的に小さくなります。これが時間方向の勾配消失です。
\end{answers}
